你把前面十六个部分读完了。合上书，有人问一句：所以这些数学到底怎么帮你做出一个靠得住的系统？

你会发现自己能讲很多具体的东西——谱定理说了什么，KKT 条件长什么样，Bellman 方程为什么是个不动点——但拼不成一个回答。这不是读得不认真。教科书必须挑一个顺序，而按学科排是唯一讲得清楚的那个；可任何一个真实的问题都不按学科发生，它横着穿过所有部分，而且穿的顺序每次都不一样。

这一部分不引入任何新数学。它做的是把同一堆内容换两种方式重新排一遍：一次按一个系统从诞生到上线的时间顺序，一次按那几个反复出现的数学动作的结构位置。两种排法各回答一个问题——``我现在卡在这一步，该翻哪一章''，以及``我读过的这些东西，其实是同一件事吗''。

\begin{center}
\begin{tikzpicture}[x=1cm,y=1cm,font=\small,>=stealth]
  % ---- 左：按时间走一遍 ----
  \foreach \i/\t in {0/定义问题,1/数据与表示,2/选模型族,3/训练,4/评估,5/部署与监控}{
    \pgfmathsetmacro{\yb}{3.62-\i*0.68}
    \draw[black!45,fill=blue!8] (0.35,\yb) rectangle (3.35,\yb+0.5);
    \node[font=\footnotesize] at (1.85,\yb+0.25) {\t};
  }
  \foreach \i in {0,1,2,3,4}{
    \pgfmathsetmacro{\ya}{3.62-\i*0.68}
    \draw[->,black!55] (1.85,\ya) -- (1.85,\ya-0.18);
  }
  \draw[->,red!65!black,thick] (0.35,0.47) -- (-0.12,0.47) -- (-0.12,3.19) -- (0.3,3.19);
  \node[rotate=90,anchor=south,font=\scriptsize,red!65!black] at (-0.28,1.8) {上线之后接上的回边};
  \node[font=\footnotesize,black!70] at (1.85,-0.35) {按一个问题走一遍};

  % ---- 右：按结构横着看 ----
  \foreach \j/\s in {0/分析,1/代数,2/概率,3/学习}{
    \pgfmathsetmacro{\xv}{7.5+\j*1.5}
    \draw[black!30,dashed] (\xv,0.35) -- (\xv,3.85);
    \node[font=\scriptsize,black!65] at (\xv,4.05) {\s};
  }
  \foreach \i/\n in {0/线性化,1/换基,2/迭代,3/不变量,4/换序}{
    \pgfmathsetmacro{\yv}{3.5-\i*0.7}
    \draw[black!70] (6.7,\yv) -- (12.6,\yv);
    \node[font=\scriptsize,anchor=east,black!85] at (6.6,\yv) {\n};
    \foreach \j in {0,1,2,3}{
      \pgfmathsetmacro{\xv}{7.5+\j*1.5}
      \fill[black!65] (\xv,\yv) circle (0.055);
    }
  }
  \node[font=\footnotesize,black!70] at (9.6,-0.35) {按一个动作横着看};
\end{tikzpicture}

\emph{同一堆内容，两种读法。左边那条路径按时间走，回答``我现在该翻哪一章''；右边那几条横线按结构走，回答``我读过的这些其实是同一件事吗''。这一部分把两种读法各走一遍。}
\end{center}

\subsection{五条贯穿始终的判断}

\textbf{拼不起来是排列方式造成的，不是理解不够。}按学科讲是最容易讲清楚的顺序，但它把一个问题的各个环节拆散到了十几个地方。一个真实的系统会在同一周之内用到凸优化的目标设计、数理统计的评估纪律、和控制论的闭环稳定性，而这三样在书里隔着八百页。换一种排列不会增加知识，但会让已有的知识变得能取用。

\textbf{症状总在病因的下游出现。}目标函数写歪了，症状要等到上线之后才显形；数据在划分之前被动过，症状出现在评估那一环；分布漂移了，症状出现在监控。所以排查故障要往回读，而不是在症状出现的那一层反复调参——离线指标掉了就换模型，往往是在给一个上游的错误做无用的补偿。第一个单元按这条线索组织，每一站都写明它做错时下游会看到什么。

\textbf{真正在起作用的数学动作只有很少几个。}在一点附近把弯的当成直的，换一组基让作用变成独立的缩放，把解写成迭代的极限，找一个不会被更新规则破坏的量，以及判断两个操作能不能交换顺序——这五个动作在十六个部分里换了名字反复出现。认出来之后，新方法就不再是新的，它多半只是同一个动作换了一件衣服。第二个单元按这条线索组织。

\textbf{每个动作都附带一张条件清单，而清单才是这本书真正教的东西。}线性化要求二阶项在你关心的尺度上可忽略，迭代要求压缩或者单调，交换次序要求质量不从尖峰或尾部逃走。学会用一个方法很便宜，知道它什么时候会\emph{安静地}失败很贵——安静是关键词，因为这些条件塌掉时程序不会报错，只会给出一个看上去正常的错答案。

\textbf{最后一条：这本书不替你做判断。}它能做的是让你知道该问什么、该去查哪一条条件。剩下那部分——这个代理指标是否忠实于真实目标，这批数据是否代表你将要面对的将来，这个近似在你的场景里够不够好——没有任何定理能回答，只能由懂业务的人来判断。把这句话放在全书最后，是因为前面所有的条件对账最终都通向它。

\subsection{两章怎么安排}

\hyperref[ch:144]{``系统的一生：症状总在病因的下游出现''}沿一条真实的工作路径走六站：把诉求写成目标函数与约束，数据与表示里被悄悄用掉的假设，选模型族其实是在选一组关于世界形状的假设，训练能保证什么不能保证什么，评估是在为一句关于未见数据的断言背书，以及上线之后系统开始改变它所预测的世界。每一站都点名判断依据在哪一章，并写出这一步做错时下游会看到的症状。最后一篇是这一章的重心，因为它引入了前五站都没覆盖的维度：部署把一个预测问题变成了反馈系统，而反馈系统不能用预测的语言分析。

\hyperref[ch:145]{``反复出现的想法：同一个动作在十六个学科里换了名字''}横着切一刀，一篇讲一个动作，展示它在互不相干的地方各叫什么名字、各自靠什么条件成立。这一章与书前的\hyperref[ch:139]{``按洞见查找：同一种结构怎样跨越学科''}分工明确：那里是查阅用的入口表，一句话加若干链接，供你在需要时定位；这里是展开来读的正文，把每个动作的失效形态和条件对账讲完整。想快速定位去前面，想真正认出来在这里。

两章顺序无所谓。手上有一个具体系统的读者从前一章进入更顺；想把读过的东西归拢成少数几件事的读者，从后一章进入更快。

读完这两章之后，可以再回头看一眼\hyperref[ch:001]{``为什么是 Beyond Formulas''}里那个最小二乘的例子。当时那句``代数上算出一条线、数值上稳定地算出它、统计上相信它能够外推，是三种不同的判断''，读起来大概只是一句提醒。走完全书再看，它已经是三个具体的部门、三套具体的条件、以及三种具体的失败方式。这本书想做的事，从头到尾就是这一件。
